\section{Grade 11: Confidence \& Interval Estimation}\label{sec:grade11-confidence}

Inferential tools for samples and populations using normal and t-student approaches in Global Economics.

Confidence intervals translate sample information into ranges of plausible population values, supporting economic and financial planning under uncertainty.

\begin{itemize}
	\item Area of simple geometric figures.
	\item Sample and population.
	\item Random variable.
	\item Continuous probability distribution.
	\item Mean and variance.
	\item Normal distribution.
	\item t-student distribution.
	\item Investment portfolios.
\end{itemize}

\begin{itemize}
	\item Point vs. interval: a single-value estimate is precise but uncertain; an interval brackets plausible values with a stated confidence level.
	\item Known variance: use Z-based critical values when population variance is available.
	\item Unknown variance: use t-student with sample standard deviation and degrees of freedom.
\end{itemize}

\subsection{Assessment Criteria}\label{sec:assessment-criteria-intervals}

\begin{enumerate}
	\item Recognises the concept of mean and standard deviation in a sample as well as in a population.
	\item Distinguishes between statistic of a population and estimating the statistic.
	\item Explains the difference between point estimation vs. interval estimation for a statistic in a population.
	\item Interprets the concept of an interval of confidence in real life applications of economy and finance.
	\item Calculates the confidence interval for the mean when the variance is known.
	\item Establishes the sample size so the error will not exceed a specified amount.
	\item Constructs the confidence interval for a proportion.
	\item Calculates confidence intervals applying t-student distribution in context situations.
	\item Estimates the confidence interval for the mean when the variance is unknown in context situations.
	\item Concludes about a statistical parameter in situations in finance and economy.
\end{enumerate}

\subsection{Support Materials}\label{sec:interval-support-materials}

\begin{itemize}
	\item \href{pages/distributions.html}{Distributions (Web Page)}: reference page with interactive probability distributions.
	\item \href{pages/confidence_interval_tool.html}{Interactive Confidence Interval Tool (Web Page)}: tool to upload data and compute confidence intervals visually.
	\item \href{pages/practice_sets.html}{Practice Sets (Web Page)}: extra practice to apply confidence interval skills and interpretation.
\end{itemize}

\vspace{6pt}
\noindent\hyperref[table-of-contents]{Back to Top}
